\section{Mina erfarenheter från LIA-perioden}
\label{sec:erfarenheter}

\subsection*{Från att lära sig verktyg till att ta ansvar}

Under LIA 1 lärde jag mig främst verktygen. Under LIA 2 fick jag i stället ansvar för att
lösningarna skulle fungera i en riktig miljö. Det gjorde stor skillnad. I skolan kan man ofta
testa sig fram utan större risk. I produktion påverkar ett fel både teamet och användarna.

Det gjorde att jag började tänka annorlunda. Jag fokuserade inte bara på att få något att fungera
just nu, utan också på att lösningen skulle vara säker, tydlig och hållbar över tid.

\subsection*{Vad som var svårast och varför}

Det svåraste under perioden var inte att skriva kod eller skapa nya resurser. Det svåraste var
att förstå \textit{varför} något inte fungerade när fel uppstod. Flera problem såg först enkla ut,
men bakom dem fanns ofta flera möjliga orsaker. Ett fel i IAM kunde till exempel se ut som ett
problem i Kubernetes, fast orsaken egentligen låg i Google Cloud.

Här lärde jag mig att felsökning måste vara metodisk. Jag behövde läsa loggar, kontrollera
behörigheter, testa en ändring i taget och sedan följa resultatet. Den arbetsmetoden gjorde mig
lugnar i pressade situationer.

\subsection*{Ett val mellan snabb lösning och rätt lösning}

I flera lägen hade jag kunnat välja en snabbare väg. Ett exempel var att ge bredare behörigheter
till service accounts för att slippa vissa fel. Det hade kanske sparat tid just då, men det hade
också gjort lösningen mindre säker och svårare att underhålla senare.

Jag valde därför att hålla fast vid en säkrare lösning med tydligare behörigheter. Det tog längre
tid, men i efterhand tycker jag att det var rätt beslut. Den erfarenheten lärde mig att en snabb
genväg ofta blir dyrare senare om den skapar teknisk skuld.

\subsection*{Konsekvenserna av driftstopp}

En annan viktig erfarenhet var att jag började se tekniska problem i ett större sammanhang.
När deployment avbröt pågående jobb var det lätt att först se det som ett rent teknikproblem.
Men ju mer jag arbetade med det, desto tydligare blev det att problemet också påverkade hur
pålitligt systemet upplevs av andra.

Om ett schemalagt jobb misslyckas kan data bli försenad eller ofullständig. Då påverkas inte
bara utvecklarna, utan även de personer som väntar på rapporter och analyser. Det gjorde att
jag såg driftstabilitet som mer än en teknisk detalj. Det handlar också om kvalitet och
förtroende.

\subsection*{Vad jag hade gjort annorlunda}

Om jag skulle göra om samma arbete hade jag tidigare byggt ett enklare end-to-end-flöde och
testat hela kedjan innan jag fördjupade mig i varje del. I början lade jag för mycket tid på att
göra vissa delar väldigt korrekta innan jag visste att helheten fungerade.

I efterhand ser jag att det hade varit bättre att först få en enkel version att fungera och sedan
förbättra den steg för steg. Det hade gjort att jag snabbare hade hittat var de största problemen
fanns.

En viktig del av denna lärdom syns också i mina pull requests. Alla gick inte hela vägen direkt.
Några stängdes och ersattes av bättre varianter. Tidigare hade jag kanske sett det som ett
misslyckande, men under LIA 2 började jag se det på ett annat sätt. En stängd pull request kan
vara ett tecken på att review-processen fungerar. Man provar en väg, får feedback, upptäcker
brister och bygger sedan en bättre lösning. Det blev särskilt tydligt i arbetet med Helm,
Workload Identity och deployment-flöden, där de första versionerna hjälpte mig att förstå vad
som verkligen behövde förbättras.

Det gjorde också att jag fick en mer realistisk bild av professionellt utvecklingsarbete. Målet
är inte att varje första försök ska vara perfekt. Målet är att ta sig fram till en lösning som är
stabil, begriplig och bra nog att användas på riktigt.

\subsection*{Kommunikation som en del av arbetet}

Jag lärde mig också att bra kommunikation är en viktig del av tekniskt arbete. Tydliga pull
requests gjorde det lättare för teamet att förstå vad jag hade ändrat, varför jag gjort det och
hur lösningen kunde testas. Det sparade tid och gjorde samarbetet bättre.

Den här erfarenheten tog jag med mig vidare under hela perioden. Jag märkte att ju tydligare
jag skrev, desto enklare blev det för andra att ge feedback och hjälpa mig framåt.